
Three of the most successful LLM formal reasoning systems use structured formal feedback from a verifier to improve performance, yet none has tested why this feedback helps. This paper introduced the oracle/regularizer distinction as a testable mechanistic framing, the locality score metric as a direct mechanistic probe, and a 3-condition 2×2 factorial design intended to test the oracle hypothesis within the BFS-Prover framework.

Attempting to execute this experiment revealed that the LeanDojo proof state extraction layer silently substituted synthetic data for real Lean4 compiler outputs. The resulting locality scores — uniformly 0.0000 across all three conditions on both evaluation datasets, with t-statistic 0.0000 and p-value 1.0000 — are a synthetic artifact. The H-E1 MUST_WORK gate failed.

The hypothesis is not falsified. The DPO training loop is implemented and validated on synthetic data. The state alignment protocol is correctly implemented. The tactic taxonomy is pre-specified. The gate evaluation logic is correct. A single environment setup task — installing the Lean4 toolchain and removing the synthetic fallback from the production code path — unblocks re-execution of H-E1.

\textbf{Summary of outputs from this run:}

\begin{enumerate}
\item The \textbf{oracle/regularizer framing} — a testable mechanistic distinction with the α-interaction prediction as a uniquely discriminating empirical test.
\end{enumerate}

\begin{enumerate}
\item The \textbf{locality score metric} — a novel mechanistic probe, not yet validated on real data.
\end{enumerate}

\begin{enumerate}
\item An \textbf{implemented and unit-tested 3-condition DPO pipeline} — ready for re-execution with real LeanDojo data.
\end{enumerate}

\begin{enumerate}
\item A \textbf{documented failure mode} — silent synthetic fallback in LLM+formal-verifier pipelines — with a proposed pre-run validation protocol.
\end{enumerate}

\textbf{Immediate future work:} Install Lean4 toolchain on the H100 cluster, verify LeanDojo produces real proof state triples, and re-execute H-E1. If H-E1 passes (LS_A > LS_P, p < 0.05), the subsequent chain (H-M1 through H-C1) tests whether the oracle mass shift translates to hard-stratum pass@1 recovery and whether the α-interaction prediction holds. Two longer-term directions are (1) testing whether the oracle effect is detectable at inference time (frozen policy, no DPO) and (2) testing cross-formalism transfer — whether the oracle mechanism holds for Verus and Dafny formalisms in Vericoding, or is Lean4-specific.
